\vspace{0.5em}\noindent\textbf{TẮT EC2 RỒI NHƯNG AWS VẪN TÍNH
PHÍ?}\par

Một hiểu lầm khá phổ biến khi mới học AWS là: \emph{EC2 đã tắt = AWS
ngừng tính phí.}

Nhưng qua những bài post và tin nhắn trong group, mình thấy rất nhiều
người vẫn nhận hóa đơn dù instance đã ở trạng thái \emph{Stopped}, thậm
chí đã \emph{Terminated}.

Có một trường hợp thực tế là một người bạn của mình đã vô tình để quên
NAT Gateway sau khi làm lab và nhận hóa đơn khoảng gần \$250.

Vấn đề nằm ở chỗ: EC2 chỉ là một phần trong hệ thống. Những tài nguyên
được tạo cùng với nó không phải lúc nào cũng tự biến mất.

\vspace{0.5em}\noindent\textbf{Stop EC2 chỉ dừng phí máy
ảo}\par

Khi EC2 ở trạng thái Stopped, AWS ngừng tính phí phần compute. Tuy
nhiên, ổ đĩa EBS vẫn được giữ lại để lần sau có thể bật máy và tiếp tục
sử dụng, vì vậy dung lượng lưu trữ vẫn bị tính phí.

Ta có thể hiểu đơn giản:

\begin{itemize}
\tightlist
\item
  Stop = Tắt máy nhưng giữ nguyên ổ cứng.
\item
  Terminate = Xóa máy, nhưng vẫn nên kiểm tra xem EBS volume có được cấu
  hình xóa cùng với instance hay không.
\end{itemize}

\vspace{0.5em}\noindent\textbf{NAT Gateway vẫn tính tiền dù EC2 đã
tắt}\par

NAT Gateway là một tài nguyên riêng trong VPC nên nó không tự dừng theo
EC2.

AWS tính phí dựa trên cả:

\begin{itemize}
\tightlist
\item
  Số giờ NAT Gateway tồn tại.
\item
  Lượng dữ liệu đi qua nó.
\end{itemize}

Ví dụ, trong bảng giá AWS tại một số region, NAT Gateway có thể tốn
khoảng \$0.045 mỗi giờ (chi phí tham khảo), tương đương hơn \$32 mỗi
tháng, ngay cả khi gần như không có traffic.

Nếu có dữ liệu đi qua, chi phí sẽ còn tăng thêm.

\vspace{0.5em}\noindent\textbf{Elastic IP cũng không miễn
phí}\par

Nếu đã cấp một Elastic IP nhưng quên release, địa chỉ đó vẫn có thể tiếp
tục tạo ra chi phí.

AWS hiện tính phí public IPv4 cả khi đang được sử dụng lẫn khi để idle.
Một địa chỉ chỉ tốn khoảng vài USD mỗi tháng, nhưng nếu nhiều địa chỉ bị
bỏ quên ở nhiều region thì chi phí sẽ cộng dồn thành một con số đáng kể.

\vspace{0.5em}\noindent\textbf{AWS Budget không phải ``hard spending
limit''}\par

AWS Budget mặc định chủ yếu dùng để gửi cảnh báo. Nó không tự động khóa
toàn bộ tài khoản khi chi phí đạt đến mức đã đặt.

Nếu muốn hệ thống tự phản ứng, chúng ta cần cấu hình thêm Budget
Actions, chẳng hạn:

\begin{itemize}
\tightlist
\item
  Stop một EC2 hoặc RDS cụ thể.
\item
  Gắn IAM policy để ngăn tạo thêm tài nguyên.
\item
  Gửi cảnh báo qua email hoặc SNS.
\end{itemize}

\vspace{0.5em}\noindent\textbf{Checklist tham khảo mình đúc kết sau mỗi buổi lab
AWS}\par

\begin{itemize}
\tightlist
\item
  Vào Cost Explorer, group theo Service và Usage Type.
\item
  Kiểm tra EC2 \(\rightarrow\) Volumes và xóa những EBS không còn cần
  thiết.
\item
  Kiểm tra VPC \(\rightarrow\) NAT Gateways.
\item
  Release các Elastic IP không còn sử dụng.
\item
  Kiểm tra đúng region vì tài nguyên ở region khác vẫn có thể tính tiền.
\item
  Bật AWS Budget và Cost Anomaly Detection ngay từ đầu.
\end{itemize}

Bài học quan trọng nhất mình rút ra là: khi xóa một service trên AWS,
đừng chỉ hỏi: \emph{``Mình đã tắt nó chưa?''} Mà hãy hỏi thêm:

\emph{``Service này đã tạo ra những tài nguyên phụ nào, và chúng có còn
đang tính phí không?''}

\texttt{\#AWS} \texttt{\#EC2} \texttt{\#CloudComputing}
\texttt{\#AWSCostOptimization}
